% Revize výukové přípravy, výchozí soubor GitHub Work-MAT, 94. hodina.
% Zdroje: Realisticky.cz 020409–020411, 020801–020802; původní test A/B koncepčně rozšířen
% Úlohy jsou vlastní / úlohy výchozí přípravy, případně nově didakticky obměněné.
\documentclass[11pt,a4paper]{article}
\input{../../sablony/pripravy/priprava.sty}
\usepackage{tikz}
\nastavPripravu{1. ročník}{Rovnice a nerovnice}{94}{Průběžný test II}
\begin{document}
\blok{Průběžný test II – absolutní hodnota, parametr a soustavy}
Uzavíráme hodiny 87--93. Zadání je v samostatných variantách A/B bez řešení; tato příprava slouží výhradně učiteli.
\blok{Organizace a bodování}
\begin{itemize}
\item 3 min zadání pokynů; 25 min samostatná práce; 12 min práce na opravách bez komentování aktuálního testu; 5 min sběr a závěr.
\item Pět úloh: dvě absolutní hodnoty (4 body), nerovnice s absolutní hodnotou (3 body), parametr (3 body), soustava (3 body), vyjádření veličiny (2 body). Celkem \textbf{15 bodů}.
\item Žáci musí uvést množinu $K$, správné intervaly, podmínku u parametru a u technického vzorce. Hodnoť i zvolenou metodu.
\end{itemize}
\blok{Ukázky typických chyb pro následný rozbor}
\begin{prikladReseniBox}
Rovnice $|x-2|+|x+2|=6$: proč nelze obě absolutní hodnoty jednoduše „zrušit“ zároveň?
\tcblower
\[\begin{aligned}
x<-2&:\quad -2x=6\Rightarrow x=-3 \\
-2\le x\le2&:\quad 4\ne6 \\
x>2&:\quad 2x=6\Rightarrow x=3 \\
K&=\{-3;3\}
\end{aligned}\]
Na prostředním intervalu jsou absolutní hodnoty vyjádřeny různými lineárními předpisy.
\end{prikladReseniBox}
\begin{zadaniReseniBox}
Diskutujte $(p-1)x=p-1$ bez dělení nulou.
\tcblower
\[\begin{aligned}
p\ne1&:\quad x=1 \\
p=1&:\quad0=0 \\
K_p&=\begin{cases}\{1\},&p\ne1,\\\mathbb R,&p=1.\end{cases}
\end{aligned}\]

\end{zadaniReseniBox}
\blok{Domácí úkol}
\begin{zadaniReseniBox}[title={Domácí příprava}]
\begin{enumerate}[label=\alph*),itemsep=2pt]
\item Vyřešte $|x+1|=3$.
\item Vyřešte $|x-1|\le2$.
\item Opravte tvrzení: $(a-3)x=0$ má pro každé $a$ pouze kořen $x=0$.
\item Vyřešte $\begin{cases}x+y=5,\\x-y=1.\end{cases}$
\item Ze vztahu $P=UI$ vyjádřete $I$ pro $U\ne0$.
\end{enumerate}
\tcblower
\begin{enumerate}[label=\alph*),itemsep=1pt]
\item
\[\begin{aligned}
x+1&=\pm3 \\
x&=-4\ \text{nebo}\ x=2 \\
K&=\{-4;2\}
\end{aligned}\]
\item
\[\begin{aligned}
-2&\le x-1\le2 \\
-1&\le x\le3 \\
K&=\langle-1;3\rangle
\end{aligned}\]
\item
\[\begin{aligned}
a\ne3&:\quad x=0 \\
a=3&:\quad 0=0 \\
K_a&=\begin{cases}\{0\},&a\ne3,\\\mathbb R,&a=3.\end{cases}
\end{aligned}\]
\item
\[\begin{aligned}
2x&=6 \\
x&=3,\quad y=2 \\
K&=\{[3;2]\}
\end{aligned}\]
\item
\[\begin{aligned}
P&=UI \\
I&=\frac PU\quad(U\ne0)
\end{aligned}\]
\end{enumerate}
\end{zadaniReseniBox}
\poznamka{Zadání ani klíč nesmí žák dostat společně. Při opravě oceň i částečné odvození po intervalech a diskusi při nulovém koeficientu. Přílohy A/B mají rozdílná čísla, ale shodné bodové váhy.}
\end{document}
